\begin{table}[tb]
\centering\footnotesize\setlength{\tabcolsep}{3pt}
\caption{全12枠の候補と送信試行．Eは完全一致，Sは認証引数除外，Oはlogin要求の由来群．}
\begin{tabular}{llll}
課題 & 条件 & 候補境界 E/S/O & 通知の送信試行 \\ \hline
写真 & N & 0/0/0 & なし \\
写真 & G & 0/0/0 & なし \\
写真 & F & 0/0/0 & なし \\
精算 & G & 0/0/0 & なし \\
精算 & F & 0/0/0 & なし \\
精算 & N & 0/0/0 & なし \\
返金 & F & 0/0/0 & なし \\
返金 & N & 0/0/0 & なし \\
返金 & G & 0/0/0 & なし \\
書出し & N & 0/1/1 & なし \\
書出し & G & 0/1/1 & 009 \\
書出し & F & 0/1/1 & 009 \\
\hline\end{tabular}
\end{table}
入力計数だけを行って予算で拒否された通知は送信試行へ数えない．試行したという記録と，推論器が返答を完了したことも区別する．観測処理の小計は保存・通知構成を含む全制御費用ではない．
